% Generated draft: requires mathematical and source review.
\documentclass[11pt]{article}
\usepackage[margin=25mm]{geometry}
\usepackage{fontspec}
\setmainfont{Latin Modern Roman}
\usepackage{amsmath,amssymb,mathtools,unicode-math}
\setmathfont{Latin Modern Math}
\usepackage{enumitem,needspace,xurl,hyperref}
\setlength{\parindent}{0pt}
\setlength{\parskip}{5pt}
\setlength{\emergencystretch}{3em}
\newcommand{\N}{\mathbb{N}}
\newcommand{\Z}{\mathbb{Z}}
\newcommand{\Q}{\mathbb{Q}}
\newcommand{\R}{\mathbb{R}}
\newcommand{\C}{\mathbb{C}}
\newcommand{\F}{\mathbb{F}}
\newcommand{\A}{\mathbb{A}}
\newcommand{\PP}{\mathbb{P}}
\newcommand{\E}{\mathbb{E}}
\newcommand{\eps}{\varepsilon}
\newcommand{\dd}{\,\mathrm{d}}
\begin{document}
\section*{UnsolvedMath ID: 700852}
\textbf{Source problem code:} catalogue entry
\subsection{Definitions and mathematical statement}
Let \(G\) be a finite group and let \(x\in G\). Define the conjugacy class \(x^G=\{g^{-1}xg:g\in G\}\). For \(a,b\in G\), let \(\langle a,b\rangle\) denote the subgroup generated by \(a\) and \(b\). The normal closure of \(x\) in \(G\) is
\[
\langle x^G\rangle=\langle g^{-1}xg:g\in G\rangle,
\]
the smallest normal subgroup of \(G\) containing \(x\).

Assume that
\[
\forall y\in x^G,\qquad |\langle x,y\rangle|\text{ is odd}.
\]
Prove, without using the classification of finite simple groups (CFSG), that
\[
|\langle x^G\rangle|\text{ is odd}.
\]
The requirement to avoid CFSG is part of the requested proof, not an additional hypothesis on \(G\).
\subsection{Short English statement}
If an element of a finite group generates a subgroup of odd order together with each of its conjugates, prove that its normal closure has odd order, without using the classification of finite simple groups.
\subsection{Sources}
\begin{itemize}
\item Supplied local UnsolvedMath registry record, permanent ID \(700852\), exact source code \texttt{catalogue entry}, titled \emph{Kourovka Notebook Problem 12.33}; upstream revision \texttt{d371ee917ac9647fbb06cd6fc529deec6f3dda82}.
\item L. S. Kazarin, Problem 12.33, \emph{The Kourovka Notebook (Unsolved Problems in Group Theory)}, 12th issue (1992). The supplied record also identifies the 21st issue edition (2026), \texttt{arXiv:1401.0300v48}. The collection description names E. I. Khukhro and V. D. Mazurov as editors.
\item The registry labels the problem \emph{open}. This is a supplied curation label, not an independently verified assessment of the implication or the availability of a proof avoiding CFSG.
\end{itemize}
\end{document}
